\section{Núcleos genealógicos de caminos: composición finita, realización unitaria y límite cilíndrico}
\label{sec:integral-genealogica-caminos-rev2}
\index{Feynman!integral genealógica de caminos}
\index{camino!categoría enriquecida}
\index{integral!cilíndrica}

La composición matricial que se demuestra en esta sección establece exactamente que cada entrada de
\(U^R\) es la suma de los pesos ordenados de las rutas de longitud \(R\).
El paso siguiente conserva la información que una ruta visible no publica y
organiza los refinamientos finitos en un espacio compacto de genealogías. De
este modo se distinguen tres resultados: la identidad combinatoria finita, la
integral cilíndrica sobre el límite de caminos y la realización de esa
integral como propagador de un sistema físico concreto.

\subsection{Categoría de rutas y peso functorial}

Sea \(\mathsf P_{\rm TPK}\) la categoría libre generada por las transiciones
admisibles del TPK. Sus objetos son estados enriquecidos de frontera y sus
morfismos son secuencias composables de transiciones. La composición es la
concatenación y el camino vacío es la identidad. Sea \(\mathsf P_X\)
la categoría de rutas visibles obtenida al aplicar el lector a los
objetos y transiciones; sus composiciones son las concatenaciones visibles.
La aplicación
\[
 q_{\rm vis}:\mathsf P_{\rm TPK}\longrightarrow\mathsf P_X
\]
publica la ruta observable y conserva en su fibra las diferencias de hoja,
orientación, acarreo, bloque, frontera y memoria.

\begin{definition}[Peso genealógico]
\label{def:peso-genealogico-caminos-rev2}
Sea \(\mathcal B\) un álgebra real normada, posiblemente no conmutativa, y
sea \(\mathbf B_{\mathcal B}\) la categoría de un solo objeto \(\ast\) cuyo
monoide de endomorfismos es
\(\operatorname{End}(\ast)=(\mathcal B,\cdot,1)\). Un peso genealógico es un funtor
\[
 W:\mathsf P_{\rm TPK}\longrightarrow\mathbf B_{\mathcal B}
\]
tal que
\[
 W(\operatorname{id}_x)=1,
 \qquad
 W(\gamma_2\circ\gamma_1)=W(\gamma_2)W(\gamma_1).
\]
El orden del producto coincide con el orden causal de la ruta.
\end{definition}

Para estados enriquecidos \(\widetilde i,\widetilde f\), el kernel a
profundidad \(R\) es
\begin{equation}
 \mathcal K_R(\widetilde f,\widetilde i)
 =\sum_{\substack{
   \gamma\in\mathsf P_{\rm TPK}(\widetilde i,\widetilde f)\\
   |\gamma|=R}}
 W(\gamma).
 \label{eq:kernel-genealogico-finito-rev2}
\end{equation}
Dos rutas con la misma palabra visible permanecen como sumandos diferentes
cuando pertenecen a fibras genealógicas distintas.

\begin{proposition}[Composición por corte enriquecido]
\label{prop:composicion-kernel-enriquecido-rev2}
Si toda ruta de longitud \(R+S\) posee un corte único después de \(R\)
transiciones, entonces
\[
 \mathcal K_{R+S}(\widetilde f,\widetilde i)
 =\sum_{\widetilde x}
 \mathcal K_S(\widetilde f,\widetilde x)
 \mathcal K_R(\widetilde x,\widetilde i).
\]
La fórmula conserva el orden causal de los factores cuando \(\mathcal B\) no
es conmutativa.
\end{proposition}

\begin{proof}
Se particiona el conjunto de rutas completas por el estado enriquecido
\(\widetilde x\) alcanzado en el corte. La factorización functorial del peso
convierte la suma sobre cada clase en el producto de los kernels parciales.
La unicidad del corte garantiza que cada ruta aparece exactamente una vez.
\end{proof}

\subsection{Construcción de la fase interna mediante el carácter real de acción sobre rutas genealógicas}

Sea \((E,J_E)\) un espacio euclídeo real provisto de un endomorfismo
ortogonal \(J_E\) tal que \(J_E^2=-I\). Si
\(\mathcal A_*:\mathsf P_{\rm TPK}\to\mathbb R\) es aditiva por
concatenación, la fase
\begin{equation}
 \Phi_J(\gamma)
 :=\exp[-J_E\mathcal A_*(\gamma)]
 =\cos\mathcal A_*(\gamma)\,I
  -\sin\mathcal A_*(\gamma)\,J_E
 \label{eq:fase-interna-caminos-rev2}
\end{equation}
es multiplicativa. Fijemos una sección positiva no nula
\(\hbar_{\rm int}\) de la línea de acción \(\mathcal L_S\)
construida anteriormente; puede utilizarse la sección
\(\hbar_{\rm ret}\) del perfil declarado.
Para una acción dimensional
\(S_{\rm int}(\gamma)=\hbar_{\rm int}\mathcal A_*(\gamma)\), esta
expresión realiza internamente la fase
\(\exp[-J_E\,S_{\rm int}(\gamma)/\hbar_{\rm int}]\). El símbolo escalar
\(\mathrm i\) aparece únicamente al elegir una carta compleja para el par
\((E,J_E)\).

\begin{proposition}[Functor de fase]
\label{prop:fase-interna-functor-rev2}
La aplicación \(\Phi_J\) satisface
\[
 \Phi_J(\operatorname{id}_x)=I,
 \qquad
 \Phi_J(\gamma_2\circ\gamma_1)
 =\Phi_J(\gamma_2)\Phi_J(\gamma_1).
\]
\end{proposition}

\begin{proof}
La acción del camino vacío es cero. Para rutas composables, la aditividad de
\(\mathcal A_*\) y la conmutación de los múltiplos de un mismo endomorfismo
\(J_E\) permiten separar la exponencial del valor sumado.
\end{proof}

\subsection{Realización unitaria sobre el registro cronológico
\texorpdfstring{\(27\times27\)}{27 x 27}}
\label{sec:realizacion-unitaria-registro-27-rev3}
\index{registro cronológico!realización unitaria}
\index{Fourier!bloques del registro cronológico}
\index{Hamiltoniano!efectivo}

Antes de pasar al límite cilíndrico, un núcleo finito admite una realización
unitaria explícita. Sea
\[
 \mathbb T_{27}^{\,2}:=(\mathbb Z/27\mathbb Z)^2,
 \qquad
 \mathcal D:=\{\mathrm N,\mathrm S,\mathrm E,\mathrm W\},
\]
y considérese el espacio real
\[
 \mathcal H_{27}^{\mathbb R,\mathrm{reg}}
 :=\ell^2(\mathbb T_{27}^{\,2};\mathbb R)
   \otimes\mathbb R^{\mathcal D}\otimes E.
\]
El endomorfismo \(J_E\) se extiende como la identidad sobre los factores de
posición y dirección, y se conserva para esa extensión la misma notación.

En la base ordenada
\((\mathrm N,\mathrm S,\mathrm E,\mathrm W)\), el mezclador direccional es
\begin{equation}
 C_4:=\frac12
 \begin{pmatrix}
  1&1&1&1\\
  1&1&-1&-1\\
  1&-1&1&-1\\
  1&-1&-1&1
 \end{pmatrix},
 \qquad C_4^{\mathsf T}C_4=I_4.
 \label{eq:mezclador-walsh-registro-27-rev3}
\end{equation}
Aquí \(C_4\) designa exclusivamente el mezclador de Walsh--Hadamard sobre la
fibra direccional; cuando actúa en
\(\mathcal H_{27}^{\mathbb R,\mathrm{reg}}\) se entiende
\(C_4\otimes I_E\).

Fíjese una carta traslacional en la que la acción de una arista elemental
\(\gamma_d\) dependa sólo de su dirección:
\[
 s_d:=S_{\rm int}(\gamma_d)\in\mathcal L_S,
 \qquad
 a_d:=\frac{s_d}{\hbar_{\rm int}}\in\mathbb R.
\]
El operador real diagonal de fase es
\[
 D_J:=\operatorname{diag}
 \left(
  e^{-J_Ea_{\mathrm N}},
  e^{-J_Ea_{\mathrm S}},
  e^{-J_Ea_{\mathrm E}},
  e^{-J_Ea_{\mathrm W}}
 \right).
\]
El desplazamiento condicionado transporta cada componente en su dirección:
\begin{align*}
 (S\psi)(x,y,\mathrm N)&=\psi(x,y-1,\mathrm N),&
 (S\psi)(x,y,\mathrm S)&=\psi(x,y+1,\mathrm S),\\
 (S\psi)(x,y,\mathrm E)&=\psi(x-1,y,\mathrm E),&
 (S\psi)(x,y,\mathrm W)&=\psi(x+1,y,\mathrm W),
\end{align*}
con todas las coordenadas tomadas módulo \(27\). El paso elemental y su
composición de doce pasos son
\begin{equation}
 U_{\rm micro}:=S D_J C_4,
 \qquad
 U_{12}^{\mathrm{reg}}:=U_{\rm micro}^{12}.
 \label{eq:paso-unitario-registro-27-rev3}
\end{equation}
El superíndice \(\mathrm{reg}\) conserva el tipo: este operador actúa sobre
el toro direccional \(27\times27\) y no se identifica con los demás objetos
dodecafásicos de la construcción ni con el endomorfismo temporal de orden \(108\).

\begin{proposition}[Unitariedad y suma exacta de caminos en el registro]
\label{prop:unitariedad-suma-registro-27-rev3}
El operador \(U_{\rm micro}\) es ortogonal sobre la realificación y unitario
en la carta compleja determinada por \(J_E\). En particular,
\(U_{12}^{\mathrm{reg}}\) es unitario. Si
\(\Xi:=\mathbb T_{27}^{\,2}\times\mathcal D\), entonces, para todo
\(n\geq0\) y \(\xi_i,\xi_f\in\Xi\),
\[
 \bigl(U_{\rm micro}^{n}\bigr)_{\xi_f,\xi_i}
 =
 \sum_{\substack{\gamma:\xi_i\to\xi_f\\|\gamma|=n}}
 W_{U_{\rm micro}}(\gamma),
\]
donde, para \(\gamma=(\xi_0,\ldots,\xi_n)\),
\[
 W_{U_{\rm micro}}(\gamma)
 :=
 (U_{\rm micro})_{\xi_n,\xi_{n-1}}
 \cdots
 (U_{\rm micro})_{\xi_1,\xi_0}.
\]
La suma recorre exactamente los caminos permitidos por los bloques no nulos
del operador.
\end{proposition}

\begin{proof}
Como \(J_E\) es ortogonal y \(J_E^2=-I\), se tiene
\(J_E^{\mathsf T}=-J_E\). Por tanto, cada bloque
\(e^{-J_Ea_d}\) es una rotación ortogonal y \(D_J\) es ortogonal.
El operador \(S\) permuta una base ortonormal y
\eqref{eq:mezclador-walsh-registro-27-rev3} demuestra que \(C_4\) es
ortogonal. En consecuencia,
\[
 U_{\rm micro}^{\mathsf T}U_{\rm micro}
 =C_4^{\mathsf T}D_J^{\mathsf T}S^{\mathsf T}
   S D_JC_4
 =I.
\]
Los tres factores conmutan con la extensión de \(J_E\); así,
\(U_{\rm micro}\) es complejo lineal en la carta
\(J_E\leftrightarrow\mathrm i\), y su ortogonalidad real equivale a
unitariedad. Toda potencia, en particular
\(U_{12}^{\mathrm{reg}}\), conserva esa propiedad.

Para \(n\geq1\), la multiplicación por bloques da
\[
 \bigl(U_{\rm micro}^{n}\bigr)_{\xi_f,\xi_i}
 =
 \sum_{\xi_1,\ldots,\xi_{n-1}}
 (U_{\rm micro})_{\xi_f,\xi_{n-1}}
 \cdots
 (U_{\rm micro})_{\xi_1,\xi_i}.
\]
Cada elección de estados intermedios determina un camino de longitud \(n\),
y cada camino determina un único producto ordenado. Para \(n=0\), el único
término es el camino vacío y su peso es la identidad.
\end{proof}

La invariancia por traslaciones descompone el problema en bloques de cuatro
componentes direccionales, tensorizados con \(E\). Para
\[
 k_x=\frac{2\pi a}{27},
 \qquad
 k_y=\frac{2\pi b}{27},
 \qquad 0\leq a,b<27,
\]
adoptamos la convención de Fourier realificada
\[
 \chi_{a,b}(x,y)
 :=\exp[-J_E(k_xx+k_yy)].
\]
La sustitución directa en el desplazamiento condicionado produce
\[
 S(k_x,k_y)
 =\operatorname{diag}
 \bigl(
  e^{J_Ek_y},e^{-J_Ek_y},
  e^{J_Ek_x},e^{-J_Ek_x}
 \bigr),
\]
y, como \(D_J\) y \(C_4\) no dependen de la posición,
\begin{equation}
 \begin{aligned}
 U_{\rm micro}(k_x,k_y)
   &=S(k_x,k_y)D_JC_4,\\
 U_{12}^{\mathrm{reg}}(k_x,k_y)
   &=U_{\rm micro}(k_x,k_y)^{12}.
 \end{aligned}
 \label{eq:bloques-fourier-registro-27-rev3}
\end{equation}
Por consiguiente,
\[
 \operatorname{Spec}U_{12}^{\mathrm{reg}}
 =
 \bigcup_{a,b=0}^{26}
 \operatorname{Spec}
 U_{12}^{\mathrm{reg}}
 \left(\frac{2\pi a}{27},\frac{2\pi b}{27}\right).
\]
La potencia matricial enumera caminos y la descomposición de Fourier publica
las cuasienergías de la misma dinámica finita.

Fijada una duración elemental \(\delta t>0\), sea
\(\Delta t_{12}:=12\,\delta t\). En cada bloque unitario elíjase una rama
espectral \(\mathfrak b\) del logaritmo. El operador
\(\operatorname{Log}_{\mathfrak b}U_{12}^{\mathrm{reg}}(k_x,k_y)\) es
anti-autoadjunto y conmuta con \(J_E\). Se define
\begin{equation}
 \begin{aligned}
 H_{\rm eff}(k_x,k_y)
 &:=
 \frac{\hbar_{\rm int}}{\Delta t_{12}}\,
 J_E\,
 \operatorname{Log}_{\mathfrak b}
 U_{12}^{\mathrm{reg}}(k_x,k_y),\\
 U_{12}^{\mathrm{reg}}(k_x,k_y)
 &=
 \exp\!\left[
  -\frac{J_E\Delta t_{12}}{\hbar_{\rm int}}\,
  H_{\rm eff}(k_x,k_y)
 \right].
 \end{aligned}
 \label{eq:hamiltoniano-efectivo-registro-27-rev3}
\end{equation}
La primera línea es autoadjunta en la carta compleja. La segunda se sigue de
\(-J_E^2=I\) y demuestra que la exponencial recupera exactamente el
propagador. La rama \(\mathfrak b\) fija la zona de cuasienergía.

Esta realización construye un propagador finito, su suma exacta de caminos y
su generador efectivo. Su identificación con un Hamiltoniano físico continuo
requiere todavía verificar las hipótesis de representación, acción y límite
que se enuncian a continuación.

\subsection{Límite cilíndrico sobre el espacio de genealogías}

Fijados los extremos, sean \(\mathcal P_n\) conjuntos finitos de caminos a
resolución \(n\), con aplicaciones de restricción sobreyectivas. El espacio
\[
 \Omega^{\rm path}_{fi}:=\varprojlim_n\mathcal P_n
\]
es compacto y totalmente disconexo. Una familia compatible de medidas
probabilísticas \(\mu_n\) sobre los cilindros determina una medida de Borel
\(\mu\) en \(\Omega^{\rm path}_{fi}\).

Sea \(\Pi_n\) la partición en cilindros de nivel \(n\). Para
\(C\in\Pi_n\), elíjanse un representante \(\gamma_C\), una acción aproximada
\(S_n(C)\) y el peso \(\mu_n(C)\). Para
\(F:\Omega^{\rm path}_{fi}\to\operatorname{End}(E)\), definimos
\begin{equation}
 \mathcal I_n(F)
 :=\sum_{C\in\Pi_n}\mu_n(C)F(\gamma_C)
 \exp\!\left[-J_E\frac{S_n(C)}{\hbar_{\rm int}}\right].
 \label{eq:integral-cilindrica-aproximante-rev2}
\end{equation}

\begin{theorem}[Integral genealógica de caminos]
\label{thm:integral-genealogica-caminos-rev2}
Supóngase que:
\begin{enumerate}
 \item las medidas cilíndricas \(\mu_n\) son proyectivamente compatibles;
 \item \(F\) es continua y acotada;
 \item existe una acción continua
 \(S:\Omega^{\rm path}_{fi}\to\mathcal L_S\) para la cual
 \[
  \max_{C\in\Pi_n}\sup_{\gamma\in C}
  \left\|\frac{S_n(C)-S(\gamma)}{\hbar_{\rm int}}\right\|
  \longrightarrow0;
 \]
 \item la malla de \(\Pi_n\) tiende a cero.
\end{enumerate}
Entonces \(\mathcal I_n(F)\) converge en norma de operador, el límite es
independiente de los representantes y define la integral de Bochner
\begin{equation}
 \int_{\Omega^{\rm path}_{fi}}
 F(\gamma)
 \exp\!\left[-J_E\frac{S(\gamma)}{\hbar_{\rm int}}\right]
 \mathrm d\mu(\gamma).
 \label{eq:integral-genealogica-caminos-rev2}
\end{equation}
\end{theorem}

\begin{proof}
La compatibilidad de las medidas finitas construye \(\mu\) sobre el álgebra
de cilindros y, por extensión, sobre la sigma-álgebra de Borel. La función
\[
 G(\gamma)=F(\gamma)
 \exp[-J_E\,S(\gamma)/\hbar_{\rm int}]
\]
es continua sobre un compacto y, por tanto, uniformemente continua y
acotada. La ortogonalidad de \(J_E\) proporciona la estimación
\[
 \|e^{-J_E a}-e^{-J_E b}\|\leq |a-b|.
\]
La aproximación uniforme de la acción y la contracción de la malla hacen
converger las sumas cilíndricas a la integral de Bochner de \(G\)
\cite{Bochner1955}.
\end{proof}

El teorema construye una integral rigurosa sobre el espacio compacto de
genealogías. Su realización como integral oscilatoria de Feynman para un
Hamiltoniano concreto se obtiene cuando los kernels finitos, la acción y la
medida de cilindros satisfacen estas hipótesis de convergencia. La expansión
matricial finita, la integral genealógica y el propagador físico quedan así
relacionados por morfismos explícitos, sin identificarse por abreviación.

\subsection{Valor publicado y memoria de la fibra}

Sea
\(\operatorname{Val}_{\rm path}:\Omega^{\rm path}_{fi}\to\mathcal Y_{\rm val}\subset\mathbb R\)
un lector de caminos Borel-medible, donde \(\mathcal Y_{\rm val}\) lleva su estructura
boreliana estándar, y sea
\(\nu=(\operatorname{Val}_{\rm path})_*\mu\). Para toda función que
factoriza por el valor, \(F=\widehat F\circ\operatorname{Val}_{\rm path}\),
se cumple
\[
 \int_{\Omega^{\rm path}_{fi}}F\,\mathrm d\mu
 =\int_{\mathcal Y_{\rm val}}\widehat F(x)\,\mathrm d\nu(x).
\]
Como \(\Omega^{\rm path}_{fi}\) es compacto metrizable y \(\mu\) es una
probabilidad de Borel, existe una familia de probabilidades condicionales
regulares \((\mu_x)_{x\in\mathcal Y_{\rm val}}\), definida para \(\nu\)-casi todo \(x\). Para
un integrando genealógico general, la desintegración de \(\mu\) adopta la forma
\begin{equation}
 \int_{\Omega^{\rm path}_{fi}}G(\gamma)\,\mathrm d\mu(\gamma)
 =\int_{\mathcal Y_{\rm val}}\left[
   \int_{\operatorname{Val}_{\rm path}^{-1}(x)}
   G(\gamma)\,\mathrm d\mu_x(\gamma)
  \right]\mathrm d\nu(x).
 \label{eq:desintegracion-genealogica-caminos-rev2}
\end{equation}
La integral interior conserva las historias distintas que publican la misma
coordenada. Este es el enlace preciso entre la suma sobre todos los caminos,
la evaluación real y la tesis según la cual el continuo es una proyección de
un dominio genealógico más rico.

\paragraph{Dominio de la integral genealógica de caminos: compacidad cilíndrica, consistencia proyectiva, convergencia de núcleos y desintegración sobre fibras}
La expansión finita, la realización unitaria sobre
\(\mathbb T_{27}^{\,2}\) y su reducción de Fourier son exactas en sus dominios
finitos. La construcción de \(H_{\rm eff}\) es
exacta una vez declaradas la estructura compleja, la acción elemental
traslacional, la duración y la rama espectral; no selecciona por sí sola un
Hamiltoniano físico continuo. La categoría enriquecida de rutas y la fase
interna preceden al ensamblaje cilíndrico. El
\cref{thm:integral-genealogica-caminos-rev2} es exacto bajo las cuatro
premisas declaradas. La realización de un modelo cuántico continuo concreto
requiere verificar esas premisas para su acción y su Hamiltoniano.
